% GENERATED FILE. Edit canonical lessons, then run npm run generate:books.
% canonical-source-hash: fnv1a64:75c33fde49b549ba
% canonical-lessons: MR-C402-ghat, MR-C402-khind, MR-C402-utar, MR-C402-nagmodi, MR-C402-sulka

\chapter{Mountain Roads}
\label{ch:mr-mountain-roads}
\hlchaptermodality{402}

\begin{chapteropening}
\textbf{By the end of this chapter:} \emph{I can talk about ghat roads, passes, slopes and a winding road.}

Complete the last lesson of chapter 402: I can talk about ghat roads, passes, slopes and a winding road.
\end{chapteropening}

\section[\texorpdfstring{\mr{घाट} (ghāṭ)}{घाट (ghāṭ)}]{\mr{घाट} (ghāṭ) — a ghat road; a mountain pass}

\label{lesson:MR-C402-ghat}



\begin{quote}
Before the new one: say the Marathi for a puncture, then the Marathi for to run out.
\end{quote}

\subsection*{You'll want to know: \mr{घाट}}
\textbf{\mr{घाट}} — \emph{ghāṭ} — \textquotedblleft{}a ghat road; a mountain pass\textquotedblright{}.

\begin{quote}
\mr{बस} \mr{हळूहळू} \mr{घाट} \mr{चढत} \mr{होती}.
\end{quote}

\begin{quote}
\emph{bas haḷūhaḷū ghāṭ caḍhat hotī.}
\end{quote}

\begin{quote}
The bus was slowly climbing the ghat.
\end{quote}

Masculine.

In the western hills a \mr{घाट} is the steep road up the slope; on a river it is a flight of steps to the water.

\begin{cousinweb}
\textbf{\mr{घाट}} (\emph{ghāṭ}) is inherited: CDIAL lists Marathi \emph{ghāṭ}, \textquotedblleft{}mountain pass, landing place\textquotedblright{}, under Sanskrit \emph{ghaṭṭa}, \textquotedblleft{}landing place, quay\textquotedblright{} (CDIAL 4414), and notes a possible Dravidian source.
\end{cousinweb}

\begin{grammarlens}[title={What you've built}]
One more word for getting around.
\end{grammarlens}

\subsection*{Your turn}
\begin{itemize}
\raggedright
  \item \emph{Say it:} \emph{ghāṭ}
  \item \emph{Say it:} \emph{ghāṭ}, once more
  \item \emph{Say it:} \emph{bas haḷūhaḷū ghāṭ caḍhat hotī.} — The bus was slowly climbing the ghat.
  \item \emph{Recall:} say the Marathi for a puncture, then the Marathi for to run out, then say \emph{ghāṭ} again
\end{itemize}

\begin{tcolorbox}[breakable,colback=teal!4,colframe=teal!35!black,title={Before you move on}]
What does \mr{घाट} mean? (\textquotedblleft{}A ghat road; a mountain pass\textquotedblright{}.) Say it once more.
\end{tcolorbox}
\section[\texorpdfstring{\mr{खिंड} (khiṇḍ)}{खिंड (khiṇḍ)}]{\mr{खिंड} (khiṇḍ) — a narrow pass; a gap}

\label{lesson:MR-C402-khind}



\begin{quote}
Before the new one: say the Marathi for to run out, then the Marathi for a ghat road.
\end{quote}

\subsection*{You'll want to know: \mr{खिंड}}
\textbf{\mr{खिंड}} — \emph{khiṇḍ} — \textquotedblleft{}a narrow pass; a gap\textquotedblright{}.

\begin{quote}
\mr{किल्ल्याच्या} \mr{मागे} \mr{एक} \mr{अरुंद} \mr{खिंड} \mr{आहे}.
\end{quote}

\begin{quote}
\emph{killyācyā māge ek arund khiṇḍ āhe.}
\end{quote}

\begin{quote}
Behind the fort there is a narrow pass.
\end{quote}

Feminine.

A gap between two hills, narrower than a \mr{घाट}.

\begin{cousinweb}
\textbf{\mr{खिंड}} (\emph{khiṇḍ}) is a Marathi word for a narrow gap between hills; its origin is not settled.
\end{cousinweb}

\begin{grammarlens}[title={What you've built}]
One more word for getting around.
\end{grammarlens}

\subsection*{Your turn}
\begin{itemize}
\raggedright
  \item \emph{Say it:} \emph{khiṇḍ}
  \item \emph{Say it:} \emph{khiṇḍ}, once more
  \item \emph{Say it:} \emph{killyācyā māge ek arund khiṇḍ āhe.} — Behind the fort there is a narrow pass.
  \item \emph{Recall:} say the Marathi for to run out, then the Marathi for a ghat road, then say \emph{khiṇḍ} again
\end{itemize}

\begin{tcolorbox}[breakable,colback=teal!4,colframe=teal!35!black,title={Before you move on}]
What does \mr{खिंड} mean? (\textquotedblleft{}A narrow pass; a gap\textquotedblright{}.) Say it once more.
\end{tcolorbox}
\section[\texorpdfstring{\mr{उतार} (utār)}{उतार (utār)}]{\mr{उतार} (utār) — a downhill slope}

\label{lesson:MR-C402-utar}



\begin{quote}
Before the new one: say the Marathi for a ghat road, then the Marathi for a narrow pass.
\end{quote}

\subsection*{You'll want to know: \mr{उतार}}
\textbf{\mr{उतार}} — \emph{utār} — \textquotedblleft{}a downhill slope\textquotedblright{}.

\begin{quote}
\mr{उतारावरून} \mr{सायकल} \mr{लवकर} \mr{खाली} \mr{येते}.
\end{quote}

\begin{quote}
\emph{utārāvarūn sāykal lavkar khālī yete.}
\end{quote}

\begin{quote}
On a downhill slope a bicycle comes down quickly.
\end{quote}

Masculine.

The opposite is \mr{चढ}, the way up; \mr{उतारावर}, on the slope.

\begin{cousinweb}
\textbf{\mr{उतार}} (\emph{utār}) is inherited: CDIAL lists Marathi \emph{utār}, \textquotedblleft{}ford, ferry fare\textquotedblright{}, under Sanskrit \emph{uttāra}, \textquotedblleft{}crossing over, landing\textquotedblright{} (CDIAL 1791); in Marathi today it is mostly a slope going down.
\end{cousinweb}

\begin{grammarlens}[title={What you've built}]
One more word for getting around.
\end{grammarlens}

\subsection*{Your turn}
\begin{itemize}
\raggedright
  \item \emph{Say it:} \emph{utār}
  \item \emph{Say it:} \emph{utār}, once more
  \item \emph{Say it:} \emph{utārāvarūn sāykal lavkar khālī yete.} — On a downhill slope a bicycle comes down quickly.
  \item \emph{Recall:} say the Marathi for a ghat road, then the Marathi for a narrow pass, then say \emph{utār} again
\end{itemize}

\begin{tcolorbox}[breakable,colback=teal!4,colframe=teal!35!black,title={Before you move on}]
What does \mr{उतार} mean? (\textquotedblleft{}A downhill slope\textquotedblright{}.) Say it once more.
\end{tcolorbox}
\section[\texorpdfstring{\mr{नागमोडी} (nāgmoḍī)}{नागमोडी (nāgmoḍī)}]{\mr{नागमोडी} (nāgmoḍī) — winding; zigzag}

\label{lesson:MR-C402-nagmodi}



\begin{quote}
Before the new one: say the Marathi for a narrow pass, then the Marathi for a downhill slope.
\end{quote}

\subsection*{You'll want to know: \mr{नागमोडी}}
\textbf{\mr{नागमोडी}} — \emph{nāgmoḍī} — \textquotedblleft{}winding; zigzag\textquotedblright{}.

\begin{quote}
\mr{घाटातला} \mr{नागमोडी} \mr{रस्ता} \mr{सुंदर} \mr{दिसतो}.
\end{quote}

\begin{quote}
\emph{ghāṭātlā nāgmoḍī rastā sundar disto.}
\end{quote}

\begin{quote}
The winding road through the ghat looks beautiful.
\end{quote}

It does not change for gender: \mr{नागमोडी} \mr{रस्ता}, \mr{नागमोडी} \mr{वाट}.

\begin{cousinweb}
\textbf{\mr{नागमोडी}} (\emph{nāgmoḍī}) is a Marathi compound: \textbf{\mr{नाग}}, \textquotedblleft{}a snake\textquotedblright{}, a tatsama from Sanskrit \emph{nāga} (Monier-Williams), and \emph{moḍī}, usually linked with \textbf{\mr{मोडणे}}, \textquotedblleft{}to bend, to break\textquotedblright{}: a road that bends like a snake. How the second part was formed is not settled.
\end{cousinweb}

\begin{grammarlens}[title={What you've built}]
One more word for getting around.
\end{grammarlens}

\subsection*{Your turn}
\begin{itemize}
\raggedright
  \item \emph{Say it:} \emph{nāgmoḍī}
  \item \emph{Say it:} \emph{nāgmoḍī}, once more
  \item \emph{Say it:} \emph{ghāṭātlā nāgmoḍī rastā sundar disto.} — The winding road through the ghat looks beautiful.
  \item \emph{Recall:} say the Marathi for a narrow pass, then the Marathi for a downhill slope, then say \emph{nāgmoḍī} again
\end{itemize}

\begin{tcolorbox}[breakable,colback=teal!4,colframe=teal!35!black,title={Before you move on}]
What does \mr{नागमोडी} mean? (\textquotedblleft{}Winding; zigzag\textquotedblright{}.) Say it once more.
\end{tcolorbox}
\section[\texorpdfstring{\mr{सुळका} (suḷkā)}{सुळका (suḷkā)}]{\mr{सुळका} (suḷkā) — a rock pinnacle}

\label{lesson:MR-C402-sulka}



\begin{quote}
Before the new one: say the Marathi for a downhill slope, then the Marathi for winding.
\end{quote}

\subsection*{You'll want to know: \mr{सुळका}}
\textbf{\mr{सुळका}} — \emph{suḷkā} — \textquotedblleft{}a rock pinnacle\textquotedblright{}.

\begin{quote}
\mr{डोंगरावर} \mr{एक} \mr{उंच} \mr{सुळका} \mr{दिसतो}.
\end{quote}

\begin{quote}
\emph{ḍoṅgarāvar ek uñc suḷkā disto.}
\end{quote}

\begin{quote}
A tall rock pinnacle stands out on the hill.
\end{quote}

Masculine.

A tall, thin peak of bare rock, loved by rock climbers.

\begin{cousinweb}
\textbf{\mr{सुळका}} (\emph{suḷkā}) is a Marathi word for a tall spike of rock; its origin is not settled.
\end{cousinweb}

\begin{grammarlens}[title={What you've built}]
One more word for getting around.
\end{grammarlens}

\subsection*{Your turn}
\begin{itemize}
\raggedright
  \item \emph{Say it:} \emph{suḷkā}
  \item \emph{Say it:} \emph{suḷkā}, once more
  \item \emph{Say it:} \emph{ḍoṅgarāvar ek uñc suḷkā disto.} — A tall rock pinnacle stands out on the hill.
  \item \emph{Recall:} say the Marathi for a downhill slope, then the Marathi for winding, then say \emph{suḷkā} again
\end{itemize}

\begin{tcolorbox}[breakable,colback=teal!4,colframe=teal!35!black,title={Before you move on}]
What does \mr{सुळका} mean? (\textquotedblleft{}A rock pinnacle\textquotedblright{}.) Say it once more.
\end{tcolorbox}
